%% ma: L10-B13-0066
%% lop: 10
%% chuong: 5
%% bai: 13
%% dang: 10.13.13
%% loai: TL
%% muc: TH
%% dap_an: "-"
%% nguon: "chuyen-de-cac-so-dac-trung-cua-mau-so-lieu-khong-ghep-nhom-toan-10-chuong-trinh-moi.pdf (GV. Phan Nhật Linh) (P03), Bài 02 Dạng 2 Bài tập tự luận 3"
%% kien_thuc: "Lớp 10 Bài 13 (số trung bình, trung vị, tứ phân vị, mốt)"
%% trang_thai: chinh-thuc
%% ghi_chu: "Tự duyệt theo yêu cầu giáo viên 22/10/2026."
\begin{ex}
Hãy tìm số trung bình, các tứ phân vị và mốt của mẫu số liệu sau:
\begin{center}
\begin{tabular}{|c|c|c|c|c|c|c|}\hline
Giá trị & $23$ & $25$ & $28$ & $31$ & $33$ & $37$\\\hline
Tần số & $6$ & $8$ & $10$ & $6$ & $4$ & $3$\\\hline
\end{tabular}
\end{center}
\loigiai{$n=37$; $\overline x=\dfrac{1047}{37}\approx28{,}3$.

Tần số tích lũy: $6$; $14$; $24$; $30$; $34$; $37$.

$Q_2=x_{19}=28$; nửa dưới gồm $18$ giá trị đầu nên $Q_1=\dfrac{x_9+x_{10}}2=25$; nửa trên gồm $18$ giá trị cuối nên $Q_3=\dfrac{x_{28}+x_{29}}2=31$.

Mốt $M_o=28$.}
\end{ex}
